% !PRINTED 610
\weberpage{610}
Den Körper $\mathfrak{K}''$ können wir erzeugen durch eine Zahl:
\weq{11}{\xi'' = \alpha\xi + \alpha'\xi',}
in der $\alpha, \alpha'$ ganze Zahlen in $\Omega$ sind.

Es sei $\mathfrak{P}''$ ein in einem $\mathfrak{p}_1$ aufgehendes Primideal in $\mathfrak{K}''$ und
$\mathfrak{P}, \mathfrak{P}'$ die durch $\mathfrak{P}''$ teilbaren Primideale in $\mathfrak{K}, \mathfrak{K}'$ und $N(\mathfrak{p}_1) = p$.
Dann ist nach 9.:
\weq{12}{\xi^p \equiv \xi \ (\mathrm{mod}\ \mathfrak{P}), \quad \xi'^p \equiv \xi' \ (\mathrm{mod}\ \mathfrak{P}'),}
und folglich ist nach (11):
\weq{13}{\xi''^p \equiv \alpha^p \xi^p + \alpha'^p \xi'^p \equiv \xi'' \ (\mathrm{mod}\ \mathfrak{P}'').}
Ist umgekehrt die Kongruenz (13) befriedigt, so folgen daraus wieder die
Kongruenzen (12), weil man $\xi$ und $\xi'$ rational [in der Form (9)] durch
$\xi''$ ausdrücken kann.

Es ist also $\mathfrak{K}''$ ebenfalls Klassenkörper, und sein Grad ist ebenso
hoch wie der von $\mathfrak{K}$ und $\mathfrak{K}'$. Es sind also $\xi$ und $\xi'$ primitive
Zahlen von $\mathfrak{K}''$, und folglich sind alle diese Körper identisch (Bd.\,I,
§ 151, 3.). Ist $\mathfrak{K}$ ein Klassenkörper, so sind die mit $\mathfrak{K}$ konjugierten
Körper auch Klassenkörper und sind daher mit $\mathfrak{K}$ identisch.

Damit ist der Satz bewiesen:

10. \gesperrt{Es gibt für eine gegebene Zahlgruppe nur einen Klassenkörper, und dies ist ein Normalkörper.}

Wir wollen nun die hauptsächlichsten Eigenschaften und Anwendungen des
Klassenkörpers in den aus der komplexen Multiplikation stammenden Fällen
näher kennen lernen.

Nehmen wir an, daß die Zahlgruppe $A$ eine andere Gruppe $A'$ als Teiler von
endlichem Index $(A, A')$ enthalte, und daß diese beiden Gruppen
Klassenkörper besitzen:
\weq{14}{\begin{aligned} \mathfrak{K} &= \mathfrak{K}(A) = \Omega(\xi), \\ \mathfrak{K}' &= \mathfrak{K}(A') = \Omega(\xi'), \end{aligned}}
dann läßt sich beweisen, daß der Körper $\mathfrak{K}$ in $\mathfrak{K}'$ enthalten ist. Sind
nämlich $\mathfrak{p}_1, \mathfrak{p}_1'$ die Primideale in $\Omega$ ersten Grades der
Hauptklassen $\overline{A}_1, \overline{A}_1'$, dann sind die $\mathfrak{p}_1'$ unter den $\mathfrak{p}_1$
enthalten. Wenn daher $p$ die Primzahl ist, in der $\mathfrak{p}_1$ und $\mathfrak{p}_1'$ aufgeht,
so ist nach 9.:
\weq{15}{\begin{aligned} \xi^p &\equiv \xi \ (\mathrm{mod}\ \mathfrak{P}) \text{ für alle } \mathfrak{p}_1 \text{ und folglich für alle } \mathfrak{p}_1', \\ \xi'^p &\equiv \xi' \ (\mathrm{mod}\ \mathfrak{P}') \text{ für alle } \mathfrak{p}_1', \end{aligned}}
wenn $\mathfrak{P}$ und $\mathfrak{P}'$ Primideale ersten Grades in $\mathfrak{K}$ und $\mathfrak{K}'$ sind.
